% Part: normal-modal-logic
% Chapter: syntax-and-semantics
% Section: relational-models

\documentclass[../../../include/open-logic-section]{subfiles}

\begin{document}

\olfileid[tr]{nml}{syn}{rel}

\olsection{Bağıntısal Modeller}

Normal kiplik mantıklarının semantiğindeki temel kavram \emph{bağıntısal
model}dir. Böyle bir model, ikili bir ``erişilebilirlik bağıntısı'' ile
birbirine bağlanan dünyalar kümesinden ve hangi !!{propositional variable}s{}nin
hangi dünyalarda ``doğru'' sayılacağını belirleyen bir atamadan oluşur.

\begin{defn}
  Temel kiplik dili için bir \emph{model}, $\mModel{M} = \tuple{W, R, V}$
  üçlüsüdür; burada
  \begin{enumerate}
  \item $W$, boş olmayan bir ``dünyalar'' kümesidir,
  \item $R$, $W$ üzerinde ikili bir erişilebilirlik bağıntısıdır ve
  \item $V$, her !!{propositional variable}~$p$'ye $V(p)$ olanaklı
    dünyalar kümesini atayan bir fonksiyondur.
  \end{enumerate}
  $Rww'$ geçerliyse $w'$ dünyasının $w$'den \emph{erişilebilir} olduğunu
  söyleriz. $w \in V(p)$ olduğunda ise $p$'nin $w$'de \emph{doğru}
  olduğunu söyleriz.
\end{defn}

Bağıntısal semantiğin büyük üstünlüğü, modellerin \olref{fig:simple}'deki
gibi basit çizimler aracılığıyla temsil edilebilmesidir. Dünyalar düğümlerle
temsil edilir; $w'$ dünyası $w$'den, ancak ve ancak $w$'den $w'$'ye bir ok
varsa erişilebilirdir. Ayrıca $w \in V(p)$ olduğunda bir düğümü (dünyayı)
$\mTrue{p}$, aksi durumda~\mFalse{p} ile etiketleriz.
\olref{fig:simple}, $W = \{w_1, w_2, w_3\}$,
$R = \{\tuple{w_1, w_2},\tuple{w_1, w_3}\}$,
$V(p) = \{w_1, w_2\}$ ve $V(q) = \{w_2\}$ olan modeli temsil eder.

\begin{figure}
  \begin{center}
    \begin{tikzpicture}[modal]
      \node[world] (w1) [label={[align=right]right:\mTrue{p}\\ \mFalse{q}}]{$w_1$};
      \node[world] (w2) [label={[align=right]right:\mTrue{p}\\ \mTrue{q}},
        above right=of w1]{$w_2$};
      \node[world] (w3) [label={[align=right]right:\mFalse{p}\\ \mFalse{q}},
        below right=of w1] {$w_3$};
      \draw[->] (w1) to (w2);
      \draw[->] (w1) to (w3);
    \end{tikzpicture}
  \end{center}
  \caption{Basit bir model.}
  \ollabel{fig:simple}
\end{figure}

\end{document}
